\section{Conclusion}

\subsection{Summary}
This project presents a comprehensive study of Transformer-based Named Entity Recognition (NER) for IT recruitment job descriptions. We construct a domain-specific dataset of 11{,}001 synthetic samples with token-level IOB annotations for five practical entity types: \textbf{ROLE}, \textbf{SKILL}, \textbf{LOC}, \textbf{EXP}, and \textbf{SALARY}. To ensure training quality, we apply a two-stage preprocessing pipeline including generation validation and tokenizer-aware label alignment for subword tokenization.

We evaluate three pre-trained Transformer backbones---\textbf{BERT-base}, \textbf{RoBERTa-base}, and \textbf{DistilBERT-base}---under six optimization settings: Base CE, Focal Loss, Dice Loss, CE + Dice, CRF, and CRF + Focal Loss. Across all architectures, the experimental results consistently show that \textbf{CE + Dice Loss} is the most effective objective. The overall best model is \textbf{RoBERTa-base + CE + Dice}, achieving an F1-score of \textbf{0.9223} with strong precision-recall balance and high token-level accuracy. BERT-base + CE + Dice follows closely (F1 = 0.9202), while DistilBERT-base + CE + Dice remains highly competitive (F1 = 0.9109) with substantially fewer parameters.

These findings demonstrate that hybrid loss design can outperform standalone objectives, and that model selection should consider both predictive quality and deployment constraints (latency, memory, and compute budget).

\subsection{Limitations}
\subsection{Future Work}